\section{Conditioned marked bridges}\label{sec:bridges}

The scalar persistence estimates must now be used with the totals that
actually occur in the arithmetic construction. Conditioning on those
totals changes the law of both ends of a path. The useful comparison is
one joint likelihood ratio, followed by integration of the remaining
bridge before taking the fractional power. This order also permits an
entire inherited measure, rather than a single entrance state, to pass
through a new clock.

\subsection{The exact binary bridge}

Fix a compact family of marks $w_+>w_->0$, their difference
$b=w_+-w_->0$, speed $c>0$, gap mean $\tau>0$, and
$\eta,\pi\in(0,1)$, with all positive quantities and complementary
probabilities bounded away from zero. Let $n\ge1$, $k=n\pi$ be an
integer, and $P=(n+1)\tau$. The clock $G$ consists of $n$ ordered
uniform points in $[0,P]$. Independently, the low marks occupy a
uniform $k$-subset of its slots. For either orientation
$\varepsilon\in\{-1,1\}$, write
\begin{equation}
 Z(t)=x+\varepsilon\{S(t)-ct\},\qquad
 y-x=\varepsilon(nw_+-bk-cP),
 \label{bridge:endpoint}
\end{equation}
where $S(t)$ is the sum of the marks at or before $t$. Both one-sided
values at a jump are included in every wall condition. Put
\begin{equation}
 \rho=\frac{(c\tau)^2}{(c\tau)^2+b^2\pi(1-\pi)},\quad
 \chi=\chi(\eta,\rho),\quad
 Q_n(x,y)=\min\left\{1,
     \left(\frac{(1+x)(1+y)}n\right)^{2\chi}\right\}.
 \label{bridge:Q}
\end{equation}
All $n+1$ gaps, including the final unmarked gap, are present.

For a deterministic grid with bounded boundary density, including
$0,P$, write $d_P(t)=\min(t,P-t)$. Given a slope $\vartheta>0$ and
$0<\zeta<1/2$, the condition $\operatorname{REG}_K$ is
\begin{equation}
 |N_a((s,t])-\lambda_a(t-s)|
 \le\vartheta(t-s)+K\{1+d_P(s)^\zeta+d_P(t)^\zeta\},
 \quad \lambda_- =k/P,\quad\lambda_+=(n-k)/P,
 \label{bridge:REG}
\end{equation}
for both colors and every grid pair $s\le t$.

\begin{theorem}[Exact binary bridges with all endpoint reserves]
\label{bridge:all-reserves}
Fix $g\ge0$ and $\zeta<\gamma<1/2$. There are fixed $h_0,K$ such
that, for every $n$, both signs, and all legal $x,y\ge h_0$,
\begin{equation}
 \mathbb E_G\mathbb P_{\rm marks}
 \{Z(t)\ge g d_P(t)^\gamma\ (0\le t\le P),
                \operatorname{REG}_K\mid G\}^{\eta}
 \asymp Q_n(x,y).
 \label{bridge:strong}
\end{equation}
The upper bound holds on dropping both the curve and regularity.
There is no upper bound on the reserves. One may replace the rates in
\eqref{bridge:REG} by prescribed rates differing from them by at most
$C/\sqrt n$, after changing the fixed constants.
\end{theorem}

We first establish the relative regularity needed in its proof.

\begin{lemma}[Relative one-end regularity]\label{bridge:relative-REG}
For the independent exponential/binary reference law with either risk
sign and bounded scaled drift on $m$ steps, set
\[
 F_m(h)=\min\{1,((1+h)/\sqrt{1+m})^{2\chi}\}.
\]
The one-end lower family from Theorem~\ref{pr:one}, with its strong
curve, compatible central separate totals, and positive inner endpoint,
can also retain two-ended local gap and color regularity, at a fixed
fraction of its $\eta$-moment, uniformly for
$h_{\rm pre}\le h\le C_{\rm res}\sqrt m$. After a fixed reserve
adjustment its curve and regularity hold in physical time.
\end{lemma}
\begin{proof}
Let $D$ read only the first $j\le m/2$ gaps and marks. Write the
positive part of the signed prefix height in unit bins $[r,r+1)$.
Discarding initial survival, the necessary remaining wall has reserve
at most $h+r+1$. The one-end upper gives
\[
 F_{m-j}(h+r+1)\le C F_m(h)(1+r)^{2\chi_{\max}},
\]
uniformly for all $h\ge h_0$. If $p_r$ is the private prefix mass
of $D$ and that bin, independence and fractional subadditivity give
\[
 \mathbb E q_D^\eta
 \le C F_m(h)\sum_{r\ge0}(1+r)^{2\chi_{\max}}\mathbb E p_r^\eta.
\]
Concavity and Cauchy--Schwarz imply
\[
 \mathbb E p_r^\eta
 \le\mathbb P(D,\text{bin }r)^\eta
 \le\mathbb P(D)^{\eta/2}
       \mathbb P(H_j^+\ge r)^{\eta/2},
\]
with the evident adjustment for $r=0$. The signed increments have
uniform exponential moments and prefix mean $O(\sqrt j)$. Their
upper tail above a fixed multiple of $\sqrt j$ is
$C\exp\{-c\min(r^2/(1+j),r)\}$. Summation proves
\begin{equation}
 \mathbb E q_D^\eta
 \le C(1+j)^C\mathbb P(D)^{\eta/2}F_m(h).
 \label{bridge:bad-prefix}
\end{equation}
This treats the unbounded positive signed heights in the negative-jump
orientation as well.

For all rank intervals $0\le i<j\le m$ require
\begin{align*}
 \left|\sum_{i<l\le j}(E_l-\tau)\right|
 &\le\epsilon_0(j-i)+K_r[1+d_m(i)^\zeta+d_m(j)^\zeta],\\
 |N_a(i,j)-\pi_a(j-i)|
 &\le\epsilon_0(j-i)+K_r[1+d_m(i)^\zeta+d_m(j)^\zeta].
\end{align*}
A fixed exponential tilt bounds the probability of any failure $D_{ij}$ by
\[
 C\exp\{-c(j-i)-cK_r[1+d_m(i)^\zeta+d_m(j)^\zeta]\}.
\]
For $j\le m/2$, use \eqref{bridge:bad-prefix}; its polynomial is
summable against the depth terms. For $i\ge m/4$, retain necessary
survival on the first $\lfloor m/8\rfloor$ steps, independent of the
failure, and apply the one-end upper. Otherwise $j-i>m/4$; drop
survival, using $F_m(h)\ge c(1+m)^{-\chi_{\max}}$ to absorb its
reciprocal and the $O(m^2)$ choices into the length tail. Fractional
subadditivity yields
\begin{equation}
 \mathbb E q_{\rm REG\,bad}^{\eta}
           \le Ce^{-cK_r}F_m(h).\label{bridge:relative-delete}
\end{equation}
This is relative to the rare survival scale. If $q$ is the retained
strong one-end mass, then
$(q-q_{\rm bad})^\eta\ge q^\eta-q_{\rm bad}^\eta$.
Choose $K_r$ after its lower constant. At least half the moment and
all its endpoint restrictions remain.

Here is the conversion to physical time. Endpoint gap inequalities
imply $i\le C(T_i+B_K)$ and
$m-i\le C(T_m-T_i+B_K)$. At a physical interval's endpoints use
the adjacent outer arrival indices. Rank depths are at most the
corresponding physical depths times a fixed constant plus $C_K$.
Applying gap regularity to the complete intervening gaps and the two
endpoint gaps gives
\[
 |\Delta t-\tau\Delta\mathrm{rank}|
 \le C\epsilon_0\Delta\mathrm{rank}
       +C_K[1+d_{T_m}(s)^\zeta+d_{T_m}(t)^\zeta].
\]
Eliminate the rank difference using $\tau\ge\tau_-$ and choose
$\epsilon_0$ small. Combining with the color inequality gives any
prescribed positive physical slope. One endpoint mark changes a count
by at most one.

For the strong curve the right endpoint of the gap following rank $j$
is at most $C(j+1)+C_K[1+(j+1)^\zeta]$. Choose its leading rank-curve
coefficient before $K_r$. The resulting physical-curve conversion
error is bounded by $\epsilon_{\rm curve}j^\gamma+C_K'$, the first
term being paid by that coefficient. Apply the rank construction
with reserve $h-C_K'$, including the fixed adverse pre-jump adjustment.
For $h\ge2C_K'+h_{\rm pre}$ its moment is comparable to $F_m(h)$;
adding $C_K'$ back to the actual risk pays the intercept. A fixed
length threshold keeps this adjustment inside the positive inner
endpoint box. No individual gap has been bounded deterministically.
\end{proof}

\begin{proof}[Proof of Theorem~\ref{bridge:all-reserves}]
Write $B(G;x)$ for the flat exact-type mass. For its upper, if both
reserves exceed $\sqrt n$, use $B\le1$. If both are at most
$2\sqrt n$, expose the first and reversed last $m=\lfloor n/8\rfloor$
marks and their distinct gaps. The two reference drifts are opposite,
with
\begin{equation}
 \delta=\varepsilon(w_+-b\pi-c\tau)
           =\frac{y-x+\varepsilon c\tau}{n},
 \label{bridge:rank-drift}
\end{equation}
and bounded scaled magnitude. The type part of the joint density
\eqref{bridge:prefix-density} bounds the private mass pointwise by
$Cq_Lq_R$. Its gap part transfers the $\eta$-moment to independent
reference ends. Theorem~\ref{pr:one} gives
\[
 \mathbb E B^\eta\le C
 \min\{1,((1+x)/\sqrt n)^{2\chi}\}
 \min\{1,((1+y)/\sqrt n)^{2\chi}\}\le C Q_n(x,y).
\]
If $x\le\sqrt n$ and $y>2\sqrt n$, then
$\delta\asymp y/n>0$, apart from bounded $n$. Expose only a prefix
of length $m=\lfloor c_0n^2/(1+y)^2\rfloor$. When this exceeds a
fixed scalar threshold, it is at most $n/4$ and
$\delta\sqrt m\le C$. Dropping all subsequent tests and using
\eqref{bridge:prefix-density} and the one-end upper gives
\[
 \mathbb E B^\eta\le C
      \min\{1,((1+x)/\sqrt m)^{2\chi}\}\le C Q_n(x,y).
\]
If $m$ is bounded, $y\gg n$, and $Q_n(x,y)\gg1$; the trivial
upper suffices. Reverse the entire bridge for the opposite asymmetric
case, reversing all $n+1$ gaps and exchanging the signs and reserves.
All remaining bounded $n$ are immediate.

For the lower, first restrict $h_0\le x,y\le C_{\rm res}\sqrt n$,
where $C_{\rm res}$ is any fixed constant. On the same two exposed
ends use Lemma~\ref{bridge:relative-REG} and the strong one-end
lower. In the positive-jump direction the rank reserve $x-w_+$
enforces every forward pre-jump value; in the reverse direction
minima occur after the negative marks. Choose the endpoint boxes
wide and compatible after $C_{\rm res}$, the laws, and the curves.
The two actual inner heights are in $[a\sqrt n,A\sqrt n]$;
separate outer gap sums and color counts are central.

For every retained complete outer state the middle has
$r=n-2m$ marks and $r+1$ gaps, a compact gap mean and interior type.
If its duration is $T'$ and its endpoints are $x_L,y_R$, its mean
pre-jump height before mark $j+1$ is exactly
\begin{equation}
 (1-j/r)x_L+(j/r)y_R
      -\varepsilon\frac{cT'(r-j)}{r(r+1)},\qquad 0\le j<r.
 \label{bridge:middle-chord}
\end{equation}
The final endpoint has height $y_R$. The correction is $O(1)$ for both
signs. Lemma~\ref{bridge:small-tube}, separately
for its $r+1$ gaps and $r$ colors, gives a fixed positive probability
of remaining above $a\sqrt n/2$, which pays the global curve
$gP^\gamma$ for large $n$.

An ordinary exact-type bridge has uniformly small regularity failure:
\begin{equation}
 \mathbb P(\operatorname{REG}_K^c)\le Ce^{-cK}.
 \label{bridge:ordinary-REG}
\end{equation}
Indeed its two colors are independent sets of uniform points with
fixed respective counts. Each interval count is binomial of mean at
most $C|I|$. Bernstein gives a bound
$C\exp\{-c|I|-cK[1+d(s)^\zeta+d(t)^\zeta]\}$ for its failure.
Bounded grid density and $\sum_{l\ge0}e^{-cl^\zeta}<\infty$
sum all pairs. Add the piece endpoints to the grid if necessary.
Choose $K$ so the failure is less than half the middle's constant
success probability. Piecewise empirical centers differ from the
full centers by $O(n^{-1/2})$; start with a smaller slope, absorb
this difference for large $n$, and split a tested interval at the
at most two junctions. Piece depths are no larger than global depths,
so the finitely many junction terms are absorbed by the intercept.

Conditional on the outer clocks put $a=q_Lq_R$, their retained
reference private mass. The complete private success mass $U$, as
a function of the middle clock, obeys
\[
 0\le U\le Ca,\qquad \mathbb E_{\rm middle}U\ge ca.
\]
The first inequality uses the all-support type density upper. The
second uses its central lower and the constant ordinary middle
success for each complete outer word. Therefore
\begin{equation}
 \mathbb E_{\rm middle}U^\eta
 \ge(Ca)^{\eta-1}\mathbb E_{\rm middle}U\ge c'a^\eta,
 \label{bridge:integrate-before-power}
\end{equation}
with zero interpreted separately. Transfer the one joint outer gap
density and average the independent reference ends. The result is
$\gg F_n(x)F_n(y)\asymp Q_n(x,y)$ in this fixed central box.

If both reserves exceed $\sqrt n$, a small tube for the whole exact
bridge remains above half its chord, pays the global curve, and has
constant ordinary probability. Delete ordinary regularity failure
using \eqref{bridge:ordinary-REG}. Its fractional moment is at least
its ordinary probability, proving the saturated case.

It remains, after reflection, that $x<\sqrt n$ and
$y>K_{\rm asym}\sqrt n$. The drift \eqref{bridge:rank-drift} is
positive and comparable to $y/n$. For small $\delta$ take
\[
 m=\lfloor c_1\delta^{-2}\rfloor,\qquad
 x'=\min(x,C_2\sqrt m),\qquad y'=y-x+x'.
\]
Lowering the entrance this way lowers the whole path. Choose
$K_{\rm asym}$ large enough that $m\le n/8$. Only the one-end
threshold $h_{\rm pre}$ is required for $x'$. The regular strong
one-end family retains
\[
 T_{\rm pre}\asymp m,\quad a\sqrt m\le H\le A\sqrt m,
 \quad |T_{\rm pre}-m\tau|+|K_{\rm pre}-m\pi|\le C\sqrt m,
\]
and has moment
\begin{equation}
 \mathbb E q_{\rm pre}^\eta
      \gg F_m(x')\asymp Q_n(x,y).\label{bridge:prefix-scale}
\end{equation}
The comparison follows from $m^{-1/2}\asymp\delta\asymp y/n$;
when entrance clipping occurs both sides are bounded below.
For each retained complete prefix set
\[
 \begin{gathered}
 n'=n-m,\qquad P'=P-T_{\rm pre},\qquad k'=k-K_{\rm pre},\\
 M'=(n'-k')w_++k'w_-,\qquad
 v'=(y'-H)/P'=\varepsilon(M'/P'-c).
 \end{gathered}
\]
The errors are $O(\delta^{-1})$, whereas $y'\asymp n\delta$.
Increasing $K_{\rm asym}$ makes $n\delta^2$ uniformly large and yields
\begin{equation}
 v'\asymp\delta,\qquad Hv'\ge a_0>0,
       \qquad P'(v')^2\ge R_{\rm tail}.
 \label{bridge:tail-data}
\end{equation}
Residual type and gap parameters stay compact. These are estimates
for every retained prefix, not assertions about a typical one.

Use the positive compact speed
$c_{\rm half}=c+\varepsilon v'/2=(c+M'/P')/2$.
On the same remaining word,
\[
 H/2+\varepsilon(S_{\rm rem}(s)-c_{\rm half}s)
 =Z_{\rm rem}(s)-(H+v's)/2.
\]
Lemma~\ref{bridge:continuation}, with entrance $H/2$, drift $v'/2$
and fixed scaled reserve $a_0/4$, gives constant ordinary probability
of $Z_{\rm rem}(s)\ge(H+v's)/2$. By \eqref{bridge:tail-data},
\[
 (H+v's)/2\ge c_*\sqrt m(1+s/m),\qquad
 d_P(T_{\rm pre}+s)^\gamma
          \le C_*m^\gamma(1+s/m)^\gamma.
\]
For $m$ above a fixed threshold this pays the entire remaining curve,
uniformly even when $n/m$ is arbitrarily large. Retain ordinary tail
regularity by \eqref{bridge:ordinary-REG}, and join it to the prefix.
Residual rates differ from full rates by
$O(\sqrt m/n+1/n)=O(n^{-1/2})$; the prefix's center conversion uses
its fixed large effective-time threshold. Apply the single prefix
density \eqref{bridge:prefix-density}. The full mass again satisfies
$0\le U\le Cq_{\rm pre}$ and
$\mathbb E_{\rm rem}U\ge cq_{\rm pre}$, where the whole exact
remaining clock and type are integrated. Equation
\eqref{bridge:integrate-before-power} and
\eqref{bridge:prefix-scale} prove the lower.

If this effective $m$ is bounded, $\delta\ge\delta_0>0$ and, for
large $n$, the physical drift $v=(y-x)/P\ge v_0>0$. Apply the
half-chord construction to the whole bridge. Choose the final $h_0$
so $h_0v_0/4\ge1$ and
\[
 h_0/2+(v_0/2)t\ge gt^\gamma\qquad(t\ge0).
\]
The finite supremum of $gt^\gamma-v_0t/2$ permits this choice.
Continuation and ordinary regularity have constant joint probability;
$Q_n(x,y)$ is also of constant order. This avoids using a rare-prefix
estimate at one step.

For clarity, first fix the regular strong one-end family, its own
threshold, and $c_1,C_2$. Its endpoint boxes fix $a_0$. Then choose
continuation constants and $K_{\rm asym}$, then the central-box
constants, then a large effective-time threshold exceeding the curve
and prefix thresholds and $(h_{\rm pre}/C_2)^2$. This fixes $v_0$.
Only now choose the final $h_0$, the ordinary middle/tail regularity
intercepts, and a large horizon threshold. For the remaining bounded
$n$, enlarge $h_0,K$ so the bounded possible excursions, curve and
counts are covered. This last enlargement changes no internal prefix
construction, which required only $h_{\rm pre}$. Finally, prescribed
rate differences $O(n^{-1/2})$ are absorbed into an initially smaller
slope at large $n$, and into $K$ at bounded $n$.
\end{proof}

\subsection{Whole inherited measures and endpoint boxes}

The next assertion concerns a measure fixed before the new clock is
read. It includes old private words and the complete terminal type
window in the same integral.

\begin{theorem}[A whole inherited measure on one clock]
\label{bridge:posterior}
Fix common $n,P,w_+,w_-,c,\varepsilon$ as above. Let $\mu$ be a
finite nonnegative measure on old states and exact types, independent
of the new clock $G$. Each state has legal endpoints satisfying
\[
 X\le x\le C_{\rm box}X,\qquad
 Y\le y\le C_{\rm box}Y,\qquad
 \epsilon\le k/n\le1-\epsilon,
\]
where $X,Y$ exceed a sufficiently large fixed threshold. Let $W=\mu(1)$
and integrate the conditional private masses over this whole measure,
before taking the power. Put $T=(1+X)(1+Y)/n$. For a sufficiently
small fixed $t_0>0$, define
\begin{equation}
 \Psi_{\rm box}=
 \begin{cases}
 1,&T\ge t_0,\\
 T^{2\chi(\eta,\rho_0)},&T<t_0,
 \end{cases}
 \quad
 q_0=\frac{w_+-cP/n}{b},\quad
 \rho_0=\frac{(c\tau)^2}{(c\tau)^2+b^2q_0(1-q_0)}.
 \label{bridge:box-factor}
\end{equation}
In the second branch with $W>0$, $q_0$ lies in a fixed compact
subinterval of $(0,1)$. The whole flat mass $q(G)$ satisfies
\begin{equation}
 \mathbb E q(G)^\eta\asymp W^\eta\Psi_{\rm box}.
 \label{bridge:whole-mass}
\end{equation}
The same comparison holds when every constituent bridge retains the
strong curve and regularity of Theorem~\ref{bridge:all-reserves}.
For $W=0$ both sides are zero.
\end{theorem}
\begin{proof}
Time reversal reduces to $\varepsilon=1$. First observe a pointwise
coordinate monotonicity on the same clock. If two lawful pairs satisfy
$x_2\ge x_1$ and $y_2\ge y_1$, couple their low subsets using the
first $k_1$ and first $k_2$ indices of one uniform permutation.
The endpoint relation gives
$b(k_2-k_1)=(x_2-x_1)-(y_2-y_1)$.
If $k_2\ge k_1$, the path difference lies between the two
nonnegative endpoint differences. If $k_2<k_1$, it is the entrance
difference plus $b$ times the number of changed low marks already
seen. Thus
\begin{equation}
 B_{k_1,x_1}(G)\le B_{k_2,x_2}(G)
 \quad\text{whenever }x_1\le x_2, y_1\le y_2.
 \label{bridge:coordinate-monotonicity}
\end{equation}
It holds at every physical time.

In the saturated branch $q\le W$. With $\nu=\mu/W$, concavity gives
\[
 \mathbb E q^\eta\ge W^\eta
          \int\mathbb E B_{k,x}^\eta\,d\nu.
\]
Each integrand is bounded below by Theorem~\ref{bridge:all-reserves},
since its endpoint product is at least $t_0$. This also proves the
strong/regular lower there.

In the rare branch the endpoint identity gives
\[
 |k/n-q_0|\le C(X+Y)/n\le C'T.
\]
Choose $t_0$ so this forces
$q_0\in[3\epsilon/4,1-3\epsilon/4]$. Set $U=nw_+-cP$ and choose
lawful enclosing bridges
\begin{align*}
 x_-&=X/2,&k_-&=\left\lfloor\frac{x_-+U-Y/2}{b}\right\rfloor,
 &y_-&=x_-+U-bk_-\in[Y/2,Y/2+b),\\
 x_+&=2C_{\rm box}X,&
 k_+&=\left\lfloor\frac{x_++U-2C_{\rm box}Y}{b}\right\rfloor,
 &y_+&=x_++U-bk_+\in[2C_{\rm box}Y,2C_{\rm box}Y+b).
\end{align*}
Choose the reserve threshold so $X/2,Y/2$ meet the bridge theorem
and $b\le Y/2$. Then $x_-\le x\le x_+$ and
$y_-\le y\le y_+$. Their type proportions differ from $q_0$ by
$O((X+Y)/n+1/n)=O(T)$; after decreasing $t_0$ and increasing the
fixed length threshold they lie in
$[\epsilon/2,1-\epsilon/2]$. The floor choices are therefore genuine
admissible types. Equation~\eqref{bridge:coordinate-monotonicity}
gives, on the one actual clock,
\[
 WB_-(G)\le q(G)\le WB_+(G).
\]
Their endpoint products are comparable to $T$. The local Lipschitz
estimate for $\chi$ gives $|\chi_\pm-\chi_0|\le CT$.
Since $T|\log T|$ is bounded, both bridge moments are comparable
to $\Psi_{\rm box}$. Taking the power proves the flat comparison.
The same exponent argument makes every constituent $Q_n(x,y)$
comparable to $\Psi_{\rm box}$. Normalized-measure concavity and
Theorem~\ref{bridge:all-reserves} give the strong/regular lower;
the upper drops these extra conditions. No monotonicity of regularity
under the subset coupling is needed.
\end{proof}

\begin{corollary}[A central terminal band]\label{bridge:terminal-band}
Let an old-state measure $\kappa$ be independent of $G$ and have mass
$m_\kappa$. In state $\omega$, let the admissible private counts form
an integer interval $J_\omega$ containing $\asymp r$ integers,
$r\ge1$, all with $|k-n\beta|\le C_{\rm type}\sqrt n$.
Assume the common endpoint boxes of Theorem~\ref{bridge:posterior}.
Give type $k$ its original mass
$p_k=\mathbb P\{\operatorname{Bin}(n,\beta)=k\}$ and let $q$ be
the entire old-state and terminal-band mass. Then, with
\[
 \rho_\beta=\frac{(c\tau)^2}
 {(c\tau)^2+b^2\beta(1-\beta)},
\]
one has
\[
 \mathbb E q^\eta\asymp m_\kappa^\eta
       \min(1,r/\sqrt n)^\eta
       \min(1,T^{2\chi(\eta,\rho_\beta)}).
\]
This also holds with the strong curve and regularity.
\end{corollary}
\begin{proof}
Uniform central Stirling gives
$\sum_{k\in J_\omega}p_k\asymp\min(1,r/\sqrt n)$ in every state.
Thus $W\asymp m_\kappa\min(1,r/\sqrt n)$. In the rare branch,
$|q_0-\beta|\le C(T+n^{-1/2})$ and $T\ge c/n$.
The exponent change times $|\log T|$ is bounded by
$C(T|\log T|+n^{-1/2}\log n)$. In the other branch the factors
are of constant order. Apply Theorem~\ref{bridge:posterior}.
\end{proof}

For clarity, the count-resolved version has a direct formulation.
On a native block of rate $\nu$, duration $L$, high mark $a$, low
mark $a-b$, and child probability $\beta$, put
$A=a-b\beta$ and $e=\nu A-c$. Given $N=n$ and entrance $x$, let
\[
 v_k=x+an-bk-cL,\qquad
 M_n(x,I)=\sum_{k:v_k\in I}\mathbb P\{\operatorname{Bin}(n,\beta)=k\}.
\]
The private terminal kernel has the exact expansion
\begin{equation}
 K_G(x,I)=\sum_{k:v_k\in I}
          \mathbb P\{\operatorname{Bin}(n,\beta)=k\}B_k(G;x).
 \label{bridge:terminal-expansion}
\end{equation}
For a nonempty band $I=[y,y+\Delta]$, with
$2b\le\Delta\le C(1+y)$, central selected types, and
$h_0\le x\le C_{\rm res}\sqrt n$ and
$h_0\le y\le y+\Delta\le C_{\rm res}\sqrt n$, monotonicity sandwiches this mass between
$M_nB_{k_{\max}}$ and $M_nB_{k_{\min}}$. Consequently
\begin{equation}
 \mathbb E_{G\mid n}K_G(x,I)^\eta
 \asymp \min(1,\Delta/\sqrt n)^\eta
       \min\{1,((1+x)(1+y)/n)^{2\chi_*}\}.
 \label{bridge:central-kernel}
\end{equation}
The exponent may be evaluated at either extreme selected type. Their
proportions differ by $O(n^{-1/2})$, so its Lipschitz change times
$\log n$ is bounded. Sufficient count-resolved conditions are
\[
 |n-\nu L|\le C_N\sqrt L,\qquad |eL|\le C_e\sqrt L,
 \qquad h_0\le x,y,\quad x,y+\Delta\le C_{\rm res}\sqrt n.
\]
Indeed
$k-\beta n=[x+A(n-\nu L)+eL-v_k]/b=O(\sqrt L)$.
Then $\tau=L/(n+1)=\nu^{-1}+O(L^{-1/2})$ and
$\pi=\beta+O(L^{-1/2})$; one may use
$\rho=A^2/(A^2+b^2\beta(1-\beta))$ in \eqref{bridge:central-kernel}.
Every selected type is central here; an uncolored central count alone
would not verify that condition.

\subsection{Allocating a fixed fraction of the endpoint budget}

\begin{corollary}[Contracted chords]\label{bridge:contracted}
On the same exact bridge write
\[
 Z(t)=\ell(t)+B(t),\quad
 \ell(t)=(1-t/P)x+(t/P)y,\quad
 B(t)=\varepsilon\{S(t)-(t/P)M\},\quad M=nw_+-bk.
\]
For $\theta\in[\theta_0,1]$, $\theta_0>0$, replacing the wall
condition in \eqref{bridge:strong} by
$\theta\ell(t)+B(t)\ge gd_P(t)^\gamma$ preserves its comparison
to the original $Q_n(x,y)$, after increasing the fixed reserve
threshold. The whole-measure version of Theorem~\ref{bridge:posterior}
also holds, with $\theta$ allowed to depend on the old state and type,
provided that measure and selection remain independent of $G$.
\end{corollary}
\begin{proof}
The exact same-word identity is
\[
 \theta\ell(t)+B(t)
 =\theta x+\varepsilon\{S(t)-c_\theta t\},\qquad
 c_\theta=\theta c+(1-\theta)M/P.
\]
This speed is a convex combination of two positive compact rates.
Theorem~\ref{bridge:all-reserves} applies with endpoints
$\theta x,\theta y$ and unchanged type, clock and regularity centers.
If $T=(1+x)(1+y)/n$ is bounded below, both its factor and the
original factor are bounded below. Otherwise
\[
 |c_\theta-c|=(1-\theta)|y-x|/P\le C(x+y)/n\le CT.
\]
Lipschitz continuity gives $|\chi_\theta-\chi|\le CT$, while
$\theta_0^2T\le(1+\theta x)(1+\theta y)/n\le T$.
Boundedness of $T|\log T|$ proves the factor comparison.
For the whole measure, the contracted event is a subset of the
original event since $\ell\ge0$, proving the upper. Normalize the
measure, apply concavity, and integrate the individual lower for the
lower. This changes an auxiliary path inequality, without changing
any original type probability or likelihood.
\end{proof}

There is also a literal complete-cell version of the flat wall and
regularity statement. Requiring them only at cell boundaries enlarges
the continuous lower event. Conversely, if widths are bounded by
$\Delta_{\max}$, a boundary wall implies continuous survival after
increasing both reserves by $c\Delta_{\max}$: use the preceding
boundary for positive jumps and the following boundary for negative
jumps. The new $Q$ is comparable to the old one. The boundary tests
read only cell/color counts; their conditional private mass depends
only on cell occupancies. Projection of uniform positions to these
occupancies therefore preserves that mass and its fractional moment
exactly.

\subsection{Ordinary bridges with physical curves}

The terminal ordinary factor has no private binary variance. We establish
its path estimate separately rather than taking a degenerate limit of
$\rho$ in a fractional theorem.

\begin{lemma}[Flat constant-mark bridge]\label{bridge:ordinary-flat}
Let $n$ uniform points lie in $[0,P]$, and let
$Z(t)=x+\varepsilon(wN(t)-ct)$ with legal endpoint
$y=x+\varepsilon(nw-cP)$. For $w,c$ in fixed positive compact ranges,
\[
 \mathbb P\{Z(t)\ge0\ (0\le t\le P)\}
 \asymp\min\{1,xy/(1+n)\},\qquad x,y\ge\max(w,1).
\]
For arbitrary nonnegative endpoints its upper holds with
$(1+x)(1+y)$ in place of $xy$.
\end{lemma}
\begin{proof}
It suffices to take $\varepsilon=1$. Scale height by $w$ and time
by $w/c$, and put $l=cP/w$, $a=x/w$, $b'=y/w=a+n-l$.
Reversing the ordered uniform points gives precisely the constraints
$\xi_j\ge(j-b')/l$. By \cite[Lemma~8.3, pp.~53--54]{Kouk}, for
integer $n\ge1$ and real $u,v\ge1$ with $u+v-n\ge1$,
\[
 n!\operatorname{vol}\{0\le\xi_1\le\cdots\le\xi_n\le1:
                      \xi_j\ge(j-u)/v\}
 \asymp\min\{1,u(u+v-n)/n\}.
\]
Use $u=b'$, $v=l$, so the other endpoint is $a$. When $l<1$
and $a,b'\ge1$, survival is automatic and the displayed factor is
bounded below. The case $n=0$ is deterministic. Raising both reserves
by $w$ proves the upper for arbitrary nonnegative endpoints, with
the same elementary small-duration cases. Reflection proves the
negative-jump assertion and reverses all $n+1$ gaps.
\end{proof}

\begin{theorem}[Ordinary strong bridge with all reserves]
\label{bridge:ordinary-strong}
Fix positive compact $w,c,\tau$, put $P=(n+1)\tau$, and take $n\ge1$
uniform points in $[0,P]$. Let $Z(t)=x+\varepsilon(wN(t)-ct)$ have
legal endpoint $y$. Fix $g>0$, $0<\zeta<\gamma<1/2$, a regularity
slope $\vartheta>0$, and a grid of bounded boundary density. There
are fixed $h_0,K$ such that, for all legal $x,y\ge h_0$ and both signs,
let $\mathcal A_{x,y}$ be the joint event
\begin{align*}
 Z(t)&\ge gd_P(t)^\gamma\qquad(0\le t\le P),\\
 |N(t)-N(s)-(n/P)(t-s)|
 &\le\vartheta(t-s)+K[1+d_P(s)^\zeta+d_P(t)^\zeta]
 \qquad\text{for all grid pairs}.
\end{align*}
Then
\begin{equation}
 \mathbb P(\mathcal A_{x,y})\asymp\beta_n(x,y):=
       \min\{1,(1+x)(1+y)/(1+n)\}.
 \label{bridge:ordinary-strong-eq}
\end{equation}
The regularity may hold at all physical pairs. Rates within
$C/\sqrt{1+n}$ of $n/P$ are also permitted. The same result holds
for a fixed contracted chord $\theta\ell+B$, with
$\theta\in[\theta_0,1]$.
\end{theorem}
\begin{proof}
The upper is Lemma~\ref{bridge:ordinary-flat}. First prove the lower
in a fixed central box $h_0\le x,y\le C_{\rm res}\sqrt n$.
Rescale to mark and speed one, changing only compact constants and
curve parameters. Write the duration as $L$ and $\lambda=n/L$.
Then $|\lambda-1|\ll n^{-1/2}$ and $L\asymp n$. The uniform
clock is a rate-$\lambda$ Poisson process conditioned on $N(L)=n$;
this single count atom has size $\asymp n^{-1/2}$.

Two elementary consequences of Lemma~\ref{bridge:ordinary-flat} are
\begin{align}
 K_\lambda(t;u,v)&\le
       \frac{C(1+u)(1+v)}{(1+t)^{3/2}},
       \label{bridge:killed}\\
 \mathbb P_u^{(1)}(\text{survival to }t)&\le
       C\min\{1,(1+u)/\sqrt{1+t}\}.
       \label{bridge:critical-one}
\end{align}
Here $K$ is a killed Poisson endpoint atom, zero unless its count is
attainable, and both orientations are allowed. For counts at least
$t/2$, multiply the conditional ballot upper by the Poisson mode
$C/\sqrt{1+t}$ and retain its denominator. Smaller counts have
an exponential Poisson lower tail; bounded $t$ is immediate. To
prove \eqref{bridge:critical-one}, sum the ballot upper at rate one:
the positive part of the terminal height has expectation at most
$u+\sqrt t+1$. For $u>\sqrt t$ use the trivial bound.

On a first-half dyadic annulus $[d,2d]$, a strong-wall failure while
surviving flat entails a first visit to height at most
$l=g(2d)^\gamma+O(1)$. At that stopped path the likelihood relative
to rate one is
\[
 \exp\{\varepsilon(Z(t)-x)\log\lambda
                 -(\lambda-1-\log\lambda)t\}.
\]
It is bounded by a constant depending on $C_{\rm res}$:
$|\log\lambda|\ll n^{-1/2}$, $x\le C_{\rm res}\sqrt n$, and
$l\ll n^\gamma$. Equation~\eqref{bridge:critical-one} until the
visit, followed by the remaining atom \eqref{bridge:killed}, and
division by the one central count atom and the flat-ballot lower,
give a relative loss at most
\begin{equation}
 C_{C_{\rm res}}(1+l)/\sqrt d.
 \label{bridge:ordinary-annulus}
\end{equation}
Explicitly, before division by the flat probability the bridge loss
is at most $C(1+x)(1+y)(1+l)/(n\sqrt d)$, and its ratio to
$\beta_n(x,y)$ is bounded in this central range.

For positive jumps no such visit occurs before $h_0/2$ once
$g(h_0/2)^\gamma\le h_0/4$. For negative jumps a visit before
$h_0/2$ requires at least $3h_0/4+(x-h_0)+O(1)$ arrivals, whereas
the mean is at most $11h_0/20$ for large $n$. Chernoff and
\eqref{bridge:killed} give relative loss
$C(1+gh_0^\gamma)e^{-ch_0}$. Reverse the argument at the right
end. The remaining annuli sum to
$C_{C_{\rm res}}(gh_0^{\gamma-1/2}+h_0^{-1/2})$.
Choose $h_0$ so the total loss is less than one quarter.

Regularity is retained relatively to this same flat event. For a
first-half interval failure $A_{s,t}$, condition on its prefix and
apply \eqref{bridge:killed} to the remaining endpoint atom. After
dividing by the count atom and flat probability,
\[
 \frac{\mathbb P_{\rm bridge}(\text{flat},A_{s,t})}
      {\mathbb P_{\rm bridge}(\text{flat})}
 \le C_{C_{\rm res}}
       \frac{\mathbb E_\lambda[(1+Z(t))_+\mathbf1_{A_{s,t}}]}{1+x}.
\]
A fixed Chernoff parameter gives
\[
 \mathbb P_\lambda(A_{s,t})
 \le2e^{-c(t-s)-cK(1+s^\zeta+t^\zeta)}.
\]
The second moment of $(1+Z(t))_+$ is at most $C(1+x+t)^2$.
Cauchy--Schwarz makes the last two displays summable over first-half
grid pairs, with total $Ce^{-cK}$. Reverse for the last half. For
crossing pairs use their exact binomial interval law: either the
length is at least $L/4$ or both endpoint depths are at least $L/4$.
Its exponential tail absorbs the $O(L^2)$ pairs and the reciprocal
flat probability, which is at most $C(1+n)$. Hence
\begin{equation}
 \mathbb P_{\rm bridge}(\text{flat},\operatorname{REG}_K^c)
 \le Ce^{-cK}\mathbb P_{\rm bridge}(\text{flat}).
 \label{bridge:ordinary-relative}
\end{equation}
Choose $K$ after the preceding strong-wall constant. This proves the
central case by relative deletion.

We need an independent-gap prefix with deterministic marks. For
$m$ independent exponential gaps, suppose
$|\varepsilon(w-c\tau)|\sqrt m\le D$ and
$h_{\rm pre}\le h\le C_{\rm pre}\sqrt m$. There is an event of
probability
\begin{equation}
 \gg\min\{1,(1+h)/\sqrt{1+m}\}\label{bridge:ordinary-prefix}
\end{equation}
on which the prefix has a strong one-ended physical curve, local
regularity, central total duration and actual endpoint in
$[a\sqrt m,A\sqrt m]$. Here is a direct construction. Conditional
on $T_m=t$, the first $m-1$ arrivals are uniform on $[0,t]$ and
the genuine $m$th mark is at the right endpoint. Choose an actual
endpoint $H$ in a sufficiently wide fixed positive $\sqrt m$ box
and set $t=(mw-\varepsilon(H-h))/c$. These durations occupy a
central Gamma window of width $\asymp\sqrt m$ and thus constant
probability by Stirling. Its constants are chosen after the reserve
and drift bounds. The inner bridge has $m-1$ marks, $m$ gaps and
pre-jump endpoint $H-\varepsilon w$. Subtract a linear height
$a'\sqrt m\,t'/t$, where $0<a'<a/2$. This changes the speed by
$\varepsilon a'\sqrt m/t$ and leaves positive central endpoints.
The central result just proved gives the order in
\eqref{bridge:ordinary-prefix}. In the first half its strong
 two-ended curve suffices; in the second half the added linear
height is at least $a'\sqrt m/2$ and dominates the one-ended curve.
The genuine final mark and its bounded endpoint adjustment are
included before fixing $h_{\rm pre}$. This proves the prefix claim
without private-mark persistence.

For the complete exact clock, exposing its first $m$ gaps has density
\[
 \frac{f_{n+1-m,\tau}(P-T_m)}{f_{n+1,\tau}(P)}.
\]
For $m\le n/8$ this has an all-support upper and, for sufficiently
large $n$, a positive lower on $|T_m-m\tau|\le C\sqrt m$.
These are the exact-reference-mean bounds of Lemma~\ref{pr:gamma};
the lower is used only above its support and length threshold.

If $x,y\ge\sqrt n$, use on the same clock the surrogate speed
$c_0=nw/P$ and endpoints $x_0=y_0=\sqrt n$.
Its centered path is $\sqrt n+\varepsilon w(N(t)-nt/P)$.
The central result has constant probability, and its event is a
subset of the desired event since the actual chord is larger.

Otherwise, outside a sufficiently large fixed central box, reverse
to $x<\sqrt n$, $y>K_{\rm asym}\sqrt n$. The reference drift
$\delta=\varepsilon(w-c\tau)=(y-x+\varepsilon c\tau)/n$
is positive. For small $\delta$ set
$m=\lfloor c_1\delta^{-2}\rfloor$ and
$x'=\min(x,C_2\sqrt m)$, $y'=y-x+x'$.
Choose $K_{\rm asym}$ so $m\le n/8$. The prefix just constructed
has probability
$\asymp\min(1,(1+x')/\sqrt m)\asymp\beta_n(x,y)$ and retains
$T_m\asymp m$, $H\asymp\sqrt m$ and a central gap sum.
For every retained prefix, the remaining duration $P'=P-T_m$ and
physical drift $v'=(y'-H)/P'$ satisfy
\[
 v'\asymp\delta,\quad Hv'\ge a_0>0,\quad P'(v')^2\ge R_{\rm tail}.
\]
The estimates follow exactly from the $O(\sqrt m)$ prefix errors
and $n\delta^2\gg1$. Apply Lemma~\ref{bridge:continuation} with
equal marks to the residual clock, speed
$(c+(n-m)w/P')/2$, entrance $H/2$ and drift $v'/2$.
It has constant probability of the half-chord bound
\[
 Z_{\rm rem}(s)\ge(H+v's)/2
      \ge c_*\sqrt m(1+s/m).
\]
Since $d_P(T_m+s)^\gamma\le C_*m^\gamma(1+s/m)^\gamma$,
a fixed large effective-time threshold pays the whole global curve.
Ordinary exact-count regularity has failure $Ce^{-cK}$, so retain
it at a constant cost. Prefix and tail center discrepancies are
$O(m^{-1/2})$ and $O(\sqrt m/n)$; reserve a small slope for them
and join the inequalities at the junction. Integrate the residual
ordinary success and then the single prefix density. This proves
the asymmetric lower above that effective-time threshold.

Below it the drift is bounded below by $\delta_0>0$, and for
sufficiently large $n$ the whole physical drift is at least $v_0>0$.
Use the same half-chord argument on the whole bridge, choosing the
final $h_0$ so
$h_0v_0/4\ge1$ and $h_0/2+v_0t/2\ge gt^\gamma$ for all $t$.
The probability and $\beta_n$ are then both of constant order.
The order of constants is: prefix boxes and thresholds; continuation
and asymmetric cutoff; effective-time threshold; then final $h_0$,
regularity intercepts and horizon threshold $n_0$. Include the
length-independent requirement
\[
 h_0\ge2+\sup_{s\ge0}(gs^\gamma-c_{\min}s).
\]
For the remaining $n<n_0$, put all points in a fixed sufficiently
short initial interval in the positive orientation. Before its end
$Z(t)\ge h_0-c_{\max}\delta_{\rm short}$; after the last point
$Z(t)=y+c(P-t)$ pays the global curve by the preceding requirement.
This event has a fixed positive probability on the bounded range.
Reflect it for negative jumps. Increase only the output regularity
intercept to make regularity automatic there. This does not feed
a reserve enlargement back into the effective-time cutoff.

To obtain all physical pairs, first use a unit mesh and sandwich an
arbitrary interval between its inward and outward mesh approximations.
Their lengths differ by at most two and endpoint depths by at most
one, so the desired inequalities follow with a fixed larger intercept
and an initially smaller slope. The permitted center change follows
as in the binary case. Finally
\[
 \theta\ell(t)+B(t)=\theta x+\varepsilon(wN(t)-c_\theta t),
 \qquad c_\theta=\theta c+(1-\theta)nw/P.
\]
This speed stays positive compact; apply the theorem with endpoints
$\theta x,\theta y$. Its ordinary factor differs only by constants.
\end{proof}

\begin{corollary}[Fixed entrance and a nonnegative terminal surplus]
\label{bridge:ordinary-surplus}
Fix $w>0$, $g\ge0$, $0<\zeta<\gamma<1/2$ and $\vartheta>0$.
There are fixed $h,K,H,\epsilon_0>0$ such that, for $n$ uniform
points on $[0,U]$ and $v=nw-U$ with $0\le v\le\epsilon_0U$,
the joint event
\begin{align*}
 h+wN(t)-t&\ge gd_U(t)^\gamma,\\
 |N(t)-N(s)-(n/U)(t-s)|
 &\le\vartheta(t-s)+K[1+d_U(s)^\zeta+d_U(t)^\zeta],\\
 |N(t)-(n/U)t|&\le H+\vartheta d_U(t)
\end{align*}
has probability $\asymp(1+v)/(1+U)$, for all sufficiently large $U$.
The same holds on a fixed complete-cell mesh of bounded width.
\end{corollary}
\begin{proof}
Use Theorem~\ref{bridge:ordinary-strong} with endpoints $h,h+v$.
First fix $h$, then choose $\epsilon_0$ so $\epsilon_0h$ is
sufficiently small. Since $n\asymp U$, its factor is
$\asymp(1+v)/(1+U)$. Take regularity slope initially below
$\vartheta/2$. Setting $s=0$ or $t=U$ and using the exact total
bounds the endpoint discrepancy by
$(\vartheta/2)d_U(t)+K'[1+d_U(t)^\zeta]$.
The inequality $K'z^\zeta\le(\vartheta/2)z+H$ gives the third
condition at no probability cost. A boundary-only wall implies
a continuous flat wall after a fixed reserve increase, which proves
the mesh upper; the continuous lower already implies the mesh lower.
\end{proof}

The all-physical regularity also yields, for $T_0=0,T_{n+1}=P$,
\[
 |T_j-T_i-\tfrac P{n+1}(j-i)|
 \le\vartheta_{\rm gap}(j-i)
       +K_{\rm gap}[1+d_n(i)^\zeta+d_n(j)^\zeta],
 \quad d_n(i)=\min(i,n+1-i).
\]
Indeed the rank-count difference at arrival endpoints is at most one.
End-interval instances of physical regularity show that time and rank
depths are comparable up to an additive constant, by absorbing
$Kz^\zeta$ into a small linear term. Solve the resulting interval
inequality; replacing $P/n$ by $P/(n+1)$ costs a fixed constant.
This includes both boundary gaps.

Optional named regularity is compatible with the ordinary lower.
Conditional on each retained uncolored regular clock, refine its slots
by a uniform exact full type in a compact simplex. A category count
in an interval is hypergeometric with mean $p_aN(I)$. Start with
uncolored slope and intercept $K_0$ fixed. Then
$N(I)\le C|I|+CK_0[1+d(s)^\zeta+d(t)^\zeta]$.
Hypergeometric Bernstein, with a named intercept $K_1$ sufficiently
larger than $K_0$, bounds each failure by
\[
 C\exp\{-c|I|-cK_1[1+d(s)^\zeta+d(t)^\zeta]\}.
\]
Summation over grid pairs makes its probability as small as desired
uniformly for every retained uncolored clock. Thus this refinement
retains a fixed fraction before any outer mixture; it does not assert
retention of an arbitrary additional private future test.

\subsection{One root count and the ordinary terminal suffix}

\begin{theorem}[Joint ordinary root and path estimate]
\label{bridge:ordinary-root}
Fix $H$, a compact product interval $[\sigma_-,\sigma_+]\subset(0,1)$,
$C_\zeta<\infty$, and $W\ge\log H$. Let the occupied alphabets be
$H\ge a_1>\cdots>a_m\ge2$, with positive lengths $L_i$, and put
\[
 \nu_i=a_i^\sigma,\quad w_i=\log a_i,\quad c_i=a_i-1,\quad
 e_i=\nu_iw_i-c_i,\quad V=\sum_iL_i,\quad D=\sum_ie_iL_i\ge0.
\]
Read independent rate-$\nu_i$ Poisson clocks in reverse block order
$m,\ldots,1$, with jumps $w_i$ and downward speed $c_i$, obtaining
$R(0)=0$ and
$X=R(V)=\sum_iw_iN_i-\sum_ic_iL_i$.
Let $E=D+\zeta$, where $|\zeta|\le C_\zeta\sqrt{1+D}$ is fixed
independently of these clocks. For a largest block $i_*$, with a fixed
tie rule, set
\[
 M=L_{i_*},\quad A=\sum_{i<i_*}L_i,\quad B=\sum_{i>i_*}L_i,
 \qquad
 \beta_D(h)=\min\left\{1,
       \frac{(1+h+B)(1+h+D+A)}{1+M}\right\}.
\]
There are fixed $V_0,h_0$ such that, uniformly for $V\ge V_0$ and
all $h\ge h_0$,
\begin{equation}
 \mathbb P\{h+R(t)\ge0\ (0\le t\le V), E-W\le X\le E\}
       \asymp (1+V)^{-1/2}\beta_D(h).
 \label{bridge:root-factor}
\end{equation}
The upper holds without the lower length threshold. The same estimate
holds conditionally on any complete external state independent of the
native clocks. That state may set $h$ using private prefix labels;
when the target is required to be uncolored, $E$ is set by its
uncolored data. The root count selected in the proof does not read $h$.
\end{theorem}
\begin{proof}
The all-count Poisson estimates are
\begin{equation}
 \sup_np_\mu(n)\ll(1+\mu)^{-1/2},\qquad
 \sup_n\frac{p_\mu(n)}{1+n}\ll(1+\mu)^{-3/2}.
 \label{bridge:poisson-modes}
\end{equation}
The first is Stirling. For the second, use $n\ge\mu/2$ to retain
its denominator, and the exponential lower tail otherwise. At a
fixed central distance from a sufficiently large mean,
$p_\mu(n)\gg(1+\mu)^{-1/2}$.

Condition on all counts except $N_*=N_{i_*}$ and write
\[
 B_{\rm act}=\sum_{i>i_*}(w_iN_i-c_iL_i),\qquad
 A_{\rm act}=-\sum_{i<i_*}(w_iN_i-c_iL_i).
\]
The largest block's actual entrance and exit reserves are
$x=h+B_{\rm act}$ and $y=h+X+A_{\rm act}$. If either is negative
there is no survival. Otherwise discard all other path tests and
use Lemma~\ref{bridge:ordinary-flat}. At most $1+W/\log2$ reservoir
counts lie in the root window. The two bounds
\eqref{bridge:poisson-modes} separately give
\[
 C(1+M)^{-1/2},\qquad
 C(1+M)^{-3/2}
 [1+h+(B_{\rm act})_+]
 [1+h+D+|\zeta|+(A_{\rm act})_+].
\]
The side sums are independent and
$\mathbb E|B_{\rm act}|\ll1+B$,
$\mathbb E|A_{\rm act}|\ll1+A$ by their Poisson means and
variances. Also $1+D+|\zeta|\ll1+D$. Average the bounds separately,
then take their minimum. Since $M\le V\le mM$, this is the upper
in \eqref{bridge:root-factor}, with constants independent of $h$.
No centrality condition has been imposed in this upper.

For the lower retain all other counts in the independent windows
\begin{equation}
 |N_i-\nu_iL_i|\le C_0\sqrt{1+L_i},\qquad i\ne i_*.
 \label{bridge:root-count-box}
\end{equation}
Choose $C_0$ so their joint mass is bounded below, using Chebyshev
and the fixed number of blocks. For each retained vector select one
count
\begin{equation}
 N_*=\left\lfloor
 \frac{\sum_ic_iL_i+D+\zeta-\sum_{i\ne i_*}w_iN_i}{w_*}
 \right\rfloor.
 \label{bridge:root-selection}
\end{equation}
Then $E-w_*<X\le E$, hence the root window holds. The exact mean
identity $\sum_iw_i\nu_iL_i=\sum_ic_iL_i+D$ gives
\[
 |N_*-\nu_*M|
 \le C\left(1+|\zeta|+\sum_{i\ne i_*}\sqrt{1+L_i}\right)
 \le C'\sqrt{1+M},
\]
since $D\ll V\ll M$. Choose $V_0$ so this is a nonnegative central
count. Its pointwise probability is $\gg(1+M)^{-1/2}$.
On the complete selected vector,
\begin{equation}
 |w_i(N_i-\nu_iL_i)|\le C_1\sqrt{1+L_i}\quad\text{for all }i.
 \label{bridge:root-complete-box}
\end{equation}
Conditioned on this vector, the original uniform positions of the
blocks are independent. Neither the windows nor the selected root
read $h$.

We next control their paths. The thresholds
$\alpha_a=\log((a-1)/\log a)/\log a$ strictly increase with $a$,
by Lemma~\ref{hier:lem:alpha}. Choose fixed disjoint small threshold
neighborhoods. At most one occupied block $i_0$ is near its threshold;
if present, one can fix $g_0>0$ so
\[
 e_i\le-g_0\ (i<i_0),\qquad e_i\ge g_0\ (i>i_0),
                  \qquad |e_{i_0}|\le g_0/8.
\]
Away from all occupied thresholds, every present drift has magnitude
at least some fixed $g_1>0$. These choices follow from finite
threshold separation and $|e_a'(\sigma)|\le H(\log H)^2$.
A single occupied alphabet is included by its one near-threshold case.

For a block with outward positive drift at least $g$, sufficiently
large $L_i$ and \eqref{bridge:root-complete-box} imply actual
endpoint increment at least $gL_i/2$. From a fixed reserve,
Lemma~\ref{bridge:ordinary-flat} gives constant survival probability,
since $N_i\ll1+L_i$. For bounded $L_i$, both selected count and
duration are bounded, so the possible loss is deterministically
bounded. Concatenating at most $H$ blocks gives an outward side path
bounded below by a fixed $-K_{\rm side}$ with constant conditional
probability. Read positive groups forward and negative groups
backward. This event is independent of the chosen $h\ge h_0$.
Since $E\ge D-C_\zeta\sqrt{1+D}\ge-C$, choosing $h_0$ once
makes both endpoint reserves sufficient.

If $i_0$ is present, write $L_0,A_0,B_0$ for its length and the
lengths on its two sides, and
\[
 A_{\rm mean}=-\sum_{i<i_0}e_iL_i,\qquad
 B_{\rm mean}=\sum_{i>i_0}e_iL_i,\qquad
 D=-A_{\rm mean}+B_{\rm mean}+e_0L_0.
\]
Here $A_{\rm mean}\asymp A_0$, $B_{\rm mean}\asymp B_0$.
Their actual values differ by $O(\sqrt{1+A_0})$ and
$O(\sqrt{1+B_0})$. The actual central junction heights are
\[
 x=h+B_{\rm act},\qquad
 y=h+D+\zeta+\epsilon_{\rm root}+A_{\rm act},
                  \quad -W\le\epsilon_{\rm root}\le0.
\]
For all $h\ge h_0$, choose $h_0$ so that
\begin{equation}
 x\asymp1+h+B_0,\qquad y\asymp1+h+D+A_0,\qquad x,y\ge\log H.
 \label{bridge:root-reserves}
\end{equation}
Indeed $g_0t-C\sqrt{1+t}\ge g_0t/2-C'$ and
$D-C_\zeta\sqrt{1+D}\ge D/2-C''$; the fixed remainders are
absorbed while retaining at least $h/2$. This is uniform in an
unbounded $h$.

The central count is at most $C(1+L_0)$. Its conditional ballot
probability is therefore at least
\[
 c\beta_0(h),\qquad
 \beta_0(h)=\min\{1,(1+h+B_0)(1+h+D+A_0)/(1+L_0)\}.
\]
At count zero the positive endpoints suffice directly. The central
bridge and two outward side events use disjoint positions after all
counts are fixed; their junctions are the actual values above.
Their joint conditional success is at least $c\beta_0(h)$.

If $i_*=i_0$, this is the desired factor. If $i_*<i_0$, then
$A_{\rm mean}\ge g_0M$ and $L_0\le M$, whence
$B_{\rm mean}\ge7g_0M/8+D$ and $B_0\gg M$.
If $i_*>i_0$, similarly
$A_{\rm mean}+D\ge7g_0M/8$ and $D+A_0\gg M$.
In either case $\beta_0(h_0)\gg1$, so monotonicity in $h$ supplies
constant success, which dominates $\beta_D(h)\le1$. These cases
also cover ties. If there is no near-threshold block, the negative
group followed in reverse and the positive group followed forward
have constant joint success. Their junction agrees because their
exact increments add to $X$. All-negative nonempty groups are
excluded by $D\ge0$; empty one-sided groups cause no difficulty.
Multiplying the conditional success by the constant outside-count
mass and the one atom in \eqref{bridge:root-selection} proves the
lower.
\end{proof}

With a fixed reserve $h=h_0$, Theorem~\ref{bridge:ordinary-root}
gives the intrinsic scales $(1+B,1+D+A)$. Its $D=\zeta=0$
specialization gives
\[
 (1+V)^{-1/2}\min\{1,(1+A)(1+B)/(1+M)\}.
\]
For a single occupied block this is $(1+V)^{-3/2}$. Thus root
centering and ordinary survival use the same count even in the
most restrictive case.

More generally, if $\mathcal F_{\rm pre}$ contains the full old
word and all its clocks, the theorem applies conditionally whenever
the native clocks are still independent with the displayed laws.
For every nonnegative $\mathcal F_{\rm pre}$-measurable weight $A_0$,
\[
 \mathbb E[A_0\mathbf1_{\{\mathrm{root, survival}\}}]
 \asymp\mathbb E[A_0(1+V)^{-1/2}\beta_D(h)].
\]
This retains the rare shared prefix measure. It neither replaces $h$
by its mean nor moves a native average through an earlier fractional
power. The lower threshold $V\ge V_0$ cannot be removed for arbitrary
allowed translates: at vanishing length a fixed negative translate
can put the root window between all attainable counts. The fixed
width $W\ge\log H$ ensures that the root selection always hits its
lattice at large length.
